\chapter{Ja oder nein, null oder mehr}

\section{Eine Entscheidung erklären}

\begin{merksatz}
\boxtitel{Geht jemand zum Facharzt?}
Die Zielgröße ist nun $y \in \{0,1\}$: ob ein Versicherter im Jahr mindestens einmal einen Facharzt aufsucht. Im Leitfall tun das $21$ Prozent.

\begin{center}
\begin{minipage}{0.9\textwidth}
\textbf{Eine Regression auf $y$ schätzt eine Wahrscheinlichkeit --- und Wahrscheinlichkeiten liegen zwischen null und eins. Die Gerade weiß das nicht.}
\end{minipage}
\end{center}
\end{merksatz}

\begin{definition}
\boxtitelm{Definition 5.1 --- Logit und Probit}
Man modelliert die Wahrscheinlichkeit als Funktion eines linearen Index:
\[ \Pr[y = 1 \mid x] = F(x^{\top}\beta), \qquad F(t) = \begin{cases} \dfrac{1}{1+e^{-t}} & \text{Logit}\\[8pt] \Phi(t) & \text{Probit}\end{cases} \]
Geschätzt wird $\beta$ durch Maximum Likelihood:
\[ \ell(\beta) = \sum_i \Bigl[y_i\ln F(x_i^{\top}\beta) + (1-y_i)\ln\bigl(1-F(x_i^{\top}\beta)\bigr)\Bigr] . \]
\end{definition}

\begin{merksatz}
\boxtitel{Die latente Deutung}
Man stelle sich eine unbeobachtete Neigung $y^{*} = x^{\top}\beta + \varepsilon$ vor; beobachtet wird nur, ob sie eine Schwelle überschreitet:
\[ y = \mathbf 1\{y^{*} > 0\} \quad\Longrightarrow\quad \Pr[y=1\mid x] = \Pr[\varepsilon > -x^{\top}\beta] = F(x^{\top}\beta) . \]

\begin{center}
{\small
\renewcommand{\arraystretch}{1.5}
\begin{tabular}{@{}P{2.4cm}P{6.2cm}@{}}
\toprule
\textbf{Probit} & $\varepsilon$ normalverteilt.\\
\textbf{Logit} & $\varepsilon$ logistisch verteilt --- fast dieselbe Form, schwerere Ränder.\\
\bottomrule
\end{tabular}}
\end{center}

\begin{center}
\begin{minipage}{0.88\textwidth}
\kursiv{\textbf{Die latente Neigung ist der Nutzen aus Band 24}: Der Versicherte geht, wenn der Nutzen des Besuchs den der Alternative übersteigt. Damit wird aus einem statistischen Modell ein Entscheidungsmodell --- und $\beta$ erhält eine Deutung als Gewicht im Nutzen. Das ist der Grund, warum die Ökonometrie diese Modelle als \enquote{diskrete Wahl} bezeichnet.}
\end{minipage}
\end{center}
\end{merksatz}

\section{Was ein Koeffizient bedeutet}

\begin{merksatz}
\boxtitel{Satz 5.2 --- der Randeffekt}
Die Änderung der Wahrscheinlichkeit je Einheit von $x_j$ beträgt
\[ \frac{\partial\Pr[y=1\mid x]}{\partial x_j} = \beta_j\,f(x^{\top}\beta) \]
mit der Dichte $f = F'$. Beim Logit ist $f(t) = F(t)\bigl(1-F(t)\bigr)$.

\schritt{Beweis}
Kettenregel. $\blacksquare$

\begin{center}
\begin{minipage}{0.9\textwidth}
\textbf{Der Koeffizient ist keine Wirkung. Die Wirkung hängt davon ab, wo man steht --- und sie ist am größten bei fünfzig Prozent.}
\end{minipage}
\end{center}

\schritt{Im Leitfall}
Logit mit Alter und Zuzahlung, $100\,000$ Versicherte:

\begin{center}
{\small
\renewcommand{\arraystretch}{1.48}
\begin{tabular}{@{}P{3.6cm}P{2.6cm}@{}}
\toprule
\textbf{Größe} & \textbf{je Euro Zuzahlung}\\
\midrule
Koeffizient $\beta$ & $-0{,}0374$\\
Randeffekt bei $p = 0{,}5$ & $-0{,}0093$\\
mittlerer Randeffekt & $\mathbf{-0{,}0059}$\\
lineares Modell & $-0{,}0059$\\
\bottomrule
\end{tabular}}
\end{center}

\begin{center}
\begin{minipage}{0.88\textwidth}
\kursiv{\textbf{Der Randeffekt am Mittel und der mittlere Randeffekt unterscheiden sich um die Hälfte} --- weil die meisten Versicherten weit unter fünfzig Prozent liegen, wo die Kurve flach ist. Und das lineare Wahrscheinlichkeitsmodell trifft den mittleren Randeffekt genau: \textbf{Für die mittlere Wirkung braucht man das Logit nicht.} Man braucht es für Vorhersagen am Rand, für Untergruppen und für die Nutzendeutung.}
\end{minipage}
\end{center}
\end{merksatz}

\section{Null oder mehr}

\begin{merksatz}
\boxtitel{Ausgaben mit Nullen}
Nun seien die Arztkosten eines Jahres die Zielgröße. Bei $32$ Prozent der Versicherten sind sie null.

\begin{center}
{\small
\renewcommand{\arraystretch}{1.48}
\begin{tabular}{@{}P{3.4cm}P{2.4cm}@{}}
\toprule
\textbf{Verfahren} & \textbf{je Euro}\\
\midrule
KQ, alle Versicherten & $-1{,}38$\\
KQ, nur Zahler & $-1{,}00$\\
\midrule
wahre Steigung der Neigung & $\mathbf{-2{,}00}$\\
\bottomrule
\end{tabular}}
\end{center}

\begin{center}
\begin{minipage}{0.9\textwidth}
\textbf{Beide Regressionen unterschätzen --- und das Weglassen der Nullen macht es schlimmer, nicht besser.}
\end{minipage}
\end{center}
\end{merksatz}

\begin{definition}
\boxtitelm{Definition 5.3 --- das Tobit-Modell}
\[ y^{*} = x^{\top}\beta + \varepsilon, \quad \varepsilon\sim\mathcal N(0,\sigma^{2}), \qquad y = \max\{y^{*},0\} \]
Die Likelihood hat zwei Arten von Beiträgen:
\[ \ell = \sum_{y_i>0}\ln\frac1\sigma\varphi\Bigl(\frac{y_i-x_i^{\top}\beta}{\sigma}\Bigr) + \sum_{y_i=0}\ln\Phi\Bigl(-\frac{x_i^{\top}\beta}{\sigma}\Bigr) . \]
\end{definition}

\begin{merksatz}
\boxtitel{Satz 5.4 --- warum die kleinsten Quadrate scheitern}
Für die Zahler gilt
\[ \mathbb{E}[y\mid x, y>0] = x^{\top}\beta + \sigma\,\lambda\Bigl(\frac{x^{\top}\beta}{\sigma}\Bigr), \qquad \lambda(t) = \frac{\varphi(t)}{\Phi(t)} . \]

\schritt{Beweis}
Für $\varepsilon\sim\mathcal N(0,\sigma^2)$ und eine Schranke $c$ gilt $\mathbb{E}[\varepsilon\mid\varepsilon > -c] = \sigma\varphi(c/\sigma)/\Phi(c/\sigma)$ --- man integriert die Dichte über den abgeschnittenen Bereich und teilt durch seine Wahrscheinlichkeit. Einsetzen mit $c = x^{\top}\beta$. $\blacksquare$

\begin{center}
\begin{minipage}{0.9\textwidth}
\textbf{Der Zusatzterm ist eine fallende Funktion von $x^{\top}\beta$ --- er wirkt wie eine ausgelassene Größe, die mit $x$ zusammenhängt.}
\end{minipage}
\end{center}

\begin{center}
\begin{minipage}{0.88\textwidth}
\kursiv{\textbf{Das ist Endogenität durch Auswahl}: Wer trotz hoher Zuzahlung noch zahlt, hat eine ungewöhnlich hohe Neigung --- ein großes $\varepsilon$. Die Auswahl der Zahler korreliert also die Störung mit dem Regressor, obwohl beide in der Gesamtbevölkerung unabhängig sind. Kapitel~6 verallgemeinert das auf Fälle, in denen die Auswahl nach einer \kursiv{anderen} Gleichung erfolgt.}
\end{minipage}
\end{center}

\schritt{Tobit im Leitfall}
\begin{center}
{\small
\renewcommand{\arraystretch}{1.48}
\begin{tabular}{@{}P{2.6cm}P{2.0cm}P{2.0cm}P{2.2cm}@{}}
\toprule
\textbf{Zuzahlung} & $\Pr[y>0]$ & $\mathbb{E}[y]$ & $\mathbb{E}[y\mid y>0]$\\
\midrule
$0$ € & $0{,}842$ & $65{,}2$ & $77{,}5$\\
$30$ € & $0{,}499$ & $23{,}9$ & $47{,}9$\\
\midrule
Steigung $\hat\beta$ & \multicolumn{3}{l}{$-2{,}01$, \ $\hat\sigma = 60{,}1$}\\
\bottomrule
\end{tabular}}
\end{center}

\begin{center}
\begin{minipage}{0.88\textwidth}
\kursiv{\textbf{Die Zuzahlung wirkt auf zwei Wegen}: Sie senkt den Anteil der Zahler von $84$ auf $50$ Prozent und die Ausgaben der verbleibenden Zahler von $78$ auf $48$ Euro. Der Tobit-Koeffizient $-2$ ist die Wirkung auf die \kursiv{Neigung}; die Wirkung auf die tatsächlichen Ausgaben beträgt $\beta\cdot\Pr[y>0]$, bei $15$ Euro Zuzahlung also $-1{,}39$ --- genau der Wert, den die naive Regression lieferte, nur mit richtiger Deutung.}
\end{minipage}
\end{center}
\end{merksatz}

\section{Der Bruchfall}

\begin{warnung}
\boxtitelr{Was die Verteilung kostet}
\begin{center}
{\small
\renewcommand{\arraystretch}{1.52}
\begin{tabular}{@{}P{3.0cm}P{5.6cm}@{}}
\toprule
\textbf{Normalverteilung} & Das Tobit-Modell ist bei anderer Störverteilung \kursiv{inkonsistent} --- anders als die kleinsten Quadrate, die nur ineffizient werden.\\
\textbf{Gleiche Streuung} & Bei Heteroskedastie ebenso. Band 17 konnte das durch robuste Standardfehler auffangen; hier geht das nicht.\\
\textbf{Ein Prozess} & Tobit unterstellt, dass derselbe Mechanismus über Null und Höhe entscheidet. Oft sind es zwei: ob man überhaupt geht, und wie viel es dann kostet.\\
\bottomrule
\end{tabular}}
\end{center}

\begin{center}
\begin{minipage}{0.88\textwidth}
\kursiv{\textbf{Die dritte Zeile ist in der Gesundheitsökonomie der wichtigste Einwand.} Ob jemand zum Arzt geht, entscheidet der Patient; wie viel es kostet, entscheidet danach weitgehend der Arzt. Zwei Entscheidungen, zwei Gleichungen --- das Zweiteilmodell, das in Aufgabe 5.3 vorgestellt wird. Es verzichtet auf die latente Neigung und gewinnt dafür Robustheit.}
\end{minipage}
\end{center}

\schritt{Und beim Logit}
\begin{center}
\begin{minipage}{0.88\textwidth}
\kursiv{Beim Logit ist das Problem subtiler: Eine ausgelassene Größe, die mit $x$ \kursiv{nicht} zusammenhängt, verzerrt die Koeffizienten trotzdem --- zur Null hin. Im Leitfall dämpft die unbeobachtete Gesundheit den Zuzahlungskoeffizienten auf $93$ Prozent des wahren Werts, obwohl sie von der Zuzahlung unabhängig ist. \textbf{Bei linearen Modellen gibt es diesen Effekt nicht}, und er macht den Vergleich von Logit-Koeffizienten zwischen Modellen mit verschiedenen Kontrollgrößen unmöglich.}
\end{minipage}
\end{center}
\end{warnung}

\section{Übungsaufgaben}

\begin{aufgabe}
\boxtitel{Aufgabe 5.1 --- Randeffekte berechnen}
Ein Logit liefert $\beta_{\text{Zuzahlung}} = -0{,}0374$. Berechnen Sie den Randeffekt bei $p = 0{,}10$, $0{,}30$ und $0{,}50$.

\loesung
\schritt{Schritt 1 --- Die Formel}
\[ \text{Randeffekt} = \beta\cdot p(1-p) \]

\schritt{Schritt 2 --- Die Werte}
\begin{center}
{\small
\renewcommand{\arraystretch}{1.48}
\begin{tabular}{@{}P{1.6cm}P{2.2cm}P{2.6cm}@{}}
\toprule
$p$ & $p(1-p)$ & \textbf{je Euro}\\
\midrule
$0{,}10$ & $0{,}090$ & $-0{,}0034$\\
$0{,}30$ & $0{,}210$ & $-0{,}0079$\\
$0{,}50$ & $0{,}250$ & $-0{,}0093$\\
\bottomrule
\end{tabular}}
\end{center}

\schritt{Schritt 3 --- Die Deutung}
\begin{center}
\begin{minipage}{0.9\textwidth}
\textbf{Dieselbe Zuzahlung wirkt bei einem Vielnutzer dreimal so stark wie bei einem Gelegenheitsnutzer.}
\end{minipage}
\end{center}

\begin{center}
\begin{minipage}{0.88\textwidth}
\kursiv{\textbf{Das ist keine empirische Erkenntnis, sondern eine Folge der Funktionsform} --- das Logit erzwingt, dass Wirkungen bei fünfzig Prozent am größten sind. Ob das stimmt, prüft man, indem man Untergruppen getrennt schätzt. Im Leitfall stimmt es: Wer gesund ist, lässt sich von zehn Euro kaum abhalten, weil er ohnehin nicht geht.}
\end{minipage}
\end{center}
\end{aufgabe}

\begin{aufgabe}
\boxtitel{Aufgabe 5.2 --- Das lineare Wahrscheinlichkeitsmodell verteidigen}
Wann darf man $y\in\{0,1\}$ einfach linear regressieren?

\loesung
\schritt{Schritt 1 --- Was das Modell leistet}
\begin{center}
{\small
\renewcommand{\arraystretch}{1.5}
\begin{tabular}{@{}P{2.8cm}P{5.8cm}@{}}
\toprule
\textbf{Mittlerer Effekt} & wird richtig geschätzt --- im Leitfall $-0{,}0059$ gegen $-0{,}0059$.\\
\textbf{Instrumente} & lassen sich ohne Umwege anwenden; beim Logit ist das heikel.\\
\textbf{Interaktionen} & haben eine klare Deutung; beim Logit nicht.\\
\bottomrule
\end{tabular}}
\end{center}

\schritt{Schritt 2 --- Was es nicht leistet}
\begin{center}
{\small
\renewcommand{\arraystretch}{1.5}
\begin{tabular}{@{}P{2.8cm}P{5.8cm}@{}}
\toprule
\textbf{Vorhersagen} & können außerhalb $[0,1]$ liegen --- im Leitfall bei $0{,}2$ Prozent der Versicherten.\\
\textbf{Ränder} & Wirkungen bei sehr kleinen oder großen $p$ sind falsch.\\
\textbf{Streuung} & Die Störung ist zwangsläufig heteroskedastisch; robuste Standardfehler sind Pflicht.\\
\bottomrule
\end{tabular}}
\end{center}

\schritt{Schritt 3 --- Die Regel}
\begin{center}
\begin{minipage}{0.9\textwidth}
\textbf{Für den mittleren kausalen Effekt: linear. Für Vorhersagen und Nutzendeutung: Logit oder Probit.}
\end{minipage}
\end{center}

\begin{center}
\begin{minipage}{0.88\textwidth}
\kursiv{\textbf{Die Ökonometrie hat sich in den letzten zwanzig Jahren dem linearen Modell zugewandt}, die Epidemiologie hält am Logit fest. Beide haben recht --- sie stellen verschiedene Fragen. Wer Odds Ratios berichtet, will Vergleiche über Studien hinweg; wer Randeffekte berichtet, will wissen, wie viele Patienten es betrifft. Band 17, Kapitel~11 hatte die Odds Ratio eingeführt; hier steht ihre Grenze.}
\end{minipage}
\end{center}
\end{aufgabe}

\begin{aufgabe}
\boxtitel{Aufgabe 5.3 --- Das Zweiteilmodell}
Formulieren Sie ein Modell, das Ob und Wieviel trennt, und vergleichen Sie es mit Tobit.

\loesung
\schritt{Schritt 1 --- Die zwei Teile}
\begin{center}
{\small
\renewcommand{\arraystretch}{1.5}
\begin{tabular}{@{}P{2.0cm}P{6.6cm}@{}}
\toprule
\textbf{Teil 1} & Logit für $\Pr[y > 0\mid x]$.\\
\textbf{Teil 2} & Regression von $\ln y$ auf $x$, nur für $y > 0$.\\
\bottomrule
\end{tabular}}
\end{center}
\[ \mathbb{E}[y\mid x] = \Pr[y>0\mid x]\cdot\mathbb{E}[y\mid x, y>0] \]

\schritt{Schritt 2 --- Der Unterschied zu Tobit}
\begin{center}
{\small
\renewcommand{\arraystretch}{1.5}
\begin{tabular}{@{}P{2.4cm}P{3.0cm}P{3.0cm}@{}}
\toprule
& \textbf{Tobit} & \textbf{Zweiteilmodell}\\
\midrule
Parameter & ein $\beta$ & zwei Sätze\\
Nullen & dieselbe Neigung & eigene Entscheidung\\
Verteilung & normal, zwingend & frei im zweiten Teil\\
Vorzeichen & gleich für beide Wege & dürfen sich unterscheiden\\
\bottomrule
\end{tabular}}
\end{center}

\schritt{Schritt 3 --- Wann welches}
\begin{center}
\begin{minipage}{0.88\textwidth}
\kursiv{\textbf{Tobit, wenn die Null eine Zensierung ist} --- ein Messgerät, das unter einer Schwelle nichts anzeigt, oder ein Wunsch, der negativ sein könnte, aber nicht beobachtet wird. \textbf{Zweiteilmodell, wenn die Null eine Entscheidung ist.} Bei Arztkosten ist sie das: Null heißt nicht \enquote{hätte gern negative Kosten gehabt}, sondern \enquote{ist nicht gegangen}. Das Krankenversicherungsexperiment wurde deshalb von Anfang an mit dem Zweiteilmodell ausgewertet --- es war dort erfunden worden.}
\end{minipage}
\end{center}

\schritt{Schritt 4 --- Die Falle im zweiten Teil}
\begin{center}
\begin{minipage}{0.88\textwidth}
\kursiv{Wer $\ln y$ regressiert, erhält $\mathbb{E}[\ln y]$, nicht $\ln\mathbb{E}[y]$ --- die Jensensche Ungleichung aus Band 25, Kapitel~12. Zurückrechnen in Euro verlangt einen Korrekturfaktor, und der hängt von der Störverteilung ab. \textbf{Die Verteilungsannahme, die man vorn hinausgeworfen hat, kehrt hinten wieder herein} --- in milderer Form, aber sie kehrt.}
\end{minipage}
\end{center}
\end{aufgabe}

\begin{geschichte}
\boxtitelg{Ein Modell und sein Name}
\Person{James Tobin} veröffentlichte das Modell aus Definition 5.3 im Jahr 1958 --- für Haushaltsausgaben auf langlebige Güter, die bei vielen Haushalten null sind.

\begin{center}
\begin{minipage}{0.9\textwidth}
\kursiv{\textbf{Den Namen \enquote{Tobit} prägte \Person{Arthur Goldberger} 1964} in Anlehnung an \enquote{Probit} --- und in Anspielung auf einen Roman von \Person{Herman Wouk}, in dem Tobin selbst unter dem Namen Tobit auftritt; beide hatten in derselben Marineschule gedient. Tobin erhielt 1981 den Nobelpreis, allerdings für anderes.}
\end{minipage}
\end{center}

\schritt{McFadden}
\begin{center}
\begin{minipage}{0.9\textwidth}
\kursiv{\Person{Daniel McFadden} gab dem Logit ab 1973 die Nutzendeutung aus Abschnitt 5.1 und erweiterte es auf viele Alternativen --- für die Vorhersage, wie viele Fahrgäste die geplante Schnellbahn von San Francisco nutzen würden. \textbf{Seine Vorhersage lag bei $6{,}3$ Prozent, die amtliche bei $15$; tatsächlich waren es $6{,}2$.} Nobelpreis 2000, gemeinsam mit \Person{Heckman}, dessen Modell das nächste Kapitel behandelt.}
\end{minipage}
\end{center}

\schritt{Und die Gesundheitsökonomie}
\begin{center}
\begin{minipage}{0.9\textwidth}
\kursiv{\Person{Naihua Duan} und Mitarbeiter entwickelten das Zweiteilmodell 1983 für das Krankenversicherungsexperiment --- weil Tobit die Nullen dort offensichtlich falsch behandelte. \textbf{Duans Korrekturfaktor für die Rücktransformation aus Aufgabe 5.3 heißt nach ihm \enquote{smearing}} und wird in der Gesundheitsökonomie bis heute benutzt.}
\end{minipage}
\end{center}
\end{geschichte}

\begin{forschung}
\boxtitelp{Diskrete Wahl heute}
\begin{center}
{\small
\renewcommand{\arraystretch}{1.45}
\begin{tabular}{@{}P{3.2cm}P{5.0cm}@{}}
\toprule
\textbf{Gemischte Logits} & Koeffizienten, die über Personen streuen --- Heterogenität der Vorlieben.\\
\textbf{Nachfrage nach Produkten} & Marktanteile als Logit-Wahrscheinlichkeiten, mit Preisen als endogenen Regressoren.\\
\textbf{Halbparametrisch} & Schätzung ohne Verteilungsannahme für $\varepsilon$.\\
\bottomrule
\end{tabular}}
\end{center}

\schritt{Zur zweiten Zeile}
\begin{center}
\begin{minipage}{0.88\textwidth}
\kursiv{Das Modell von \Person{Berry}, \Person{Levinsohn} und \Person{Pakes} aus dem Jahr 1995 verbindet die diskrete Wahl mit den Instrumenten aus Kapitel~3: Konsumenten wählen nach Logit, Preise sind endogen, Instrumente sind die Eigenschaften konkurrierender Produkte. \textbf{Es ist das Standardwerkzeug der empirischen Industrieökonomik} und wurde zur Bewertung von Fusionen, Zöllen und Arzneimittelpreisen eingesetzt --- jedes Mal mit allen Annahmen dieses und des dritten Kapitels zugleich.}
\end{minipage}
\end{center}
\end{forschung}

\begin{merksatz}
\boxtitel{Was Sie jetzt können}
\begin{center}
\renewcommand{\arraystretch}{1.5}
\begin{tabular}{@{}P{3.0cm}P{5.6cm}@{}}
\toprule
\textbf{Definition 5.1} & Logit, Probit, latente Neigung\\
\textbf{Satz 5.2} & Koeffizient ist keine Wirkung\\
\textbf{Definition 5.3} & Tobit: Zensierung bei null\\
\textbf{Satz 5.4} & Auswahl erzeugt Endogenität\\
\textbf{Der Bruchfall} & Verteilung kostet Robustheit\\
\bottomrule
\end{tabular}
\end{center}

\begin{center}
\begin{minipage}{0.88\textwidth}
\kursiv{Satz 5.4 zeigte, dass die Auswahl der Beobachtungen eine Störgröße erzeugt. Das nächste Kapitel behandelt den Fall, dass die Auswahl nach einer eigenen Regel erfolgt --- wer in der Stichprobe ist, hat sich dafür entschieden.}
\end{minipage}
\end{center}
\end{merksatz}
